\documentclass[11pt]{article}

\usepackage[T1]{fontenc}
\usepackage[utf8]{inputenc}
\usepackage{lmodern}
\usepackage[margin=1.05in]{geometry}
\usepackage{microtype}
\usepackage{amsmath,amssymb,amsthm,mathtools}
\usepackage{enumitem}
\usepackage{xcolor}
\usepackage[hidelinks]{hyperref}

\hypersetup{pdftitle={Connectivity threshold for G(n,p)},pdfauthor={ChatGPT},pdfsubject={Self-contained proof}}

\setlength{\parskip}{0.35em}
\setlength{\parindent}{0pt}
\setlist[enumerate]{itemsep=0.25em,topsep=0.35em,leftmargin=2.2em}

\newcommand{\PP}{\mathbb{P}}
\newcommand{\EE}{\mathbb{E}}
\newcommand{\Var}{\operatorname{Var}}
\newcommand{\one}{\mathbf{1}}
\newcommand{\bbinom}[2]{\binom{#1}{#2}}

\theoremstyle{plain}
\newtheorem{theorem}{Theorem}
\newtheorem{lemma}[theorem]{Lemma}
\newtheorem{proposition}[theorem]{Proposition}
\newtheorem{corollary}[theorem]{Corollary}

\theoremstyle{definition}
\newtheorem{remark}[theorem]{Remark}

\title{Connectivity threshold for \texorpdfstring{$G_{n,p}$}{G(n,p)}\\[0.35em]
\large A self-contained proof}
\author{}
\date{March 2026}

\begin{document}
\maketitle

\begin{abstract}
We prove the classical connectivity threshold for the Erd\H{o}s-R\'enyi random graph in the form
\[
\PP\bigl(G(n,p_n)\text{ is connected}\bigr)
\to
\begin{cases}
0, & np_n-\log n\to -\infty,\\
\exp(-e^{-c}), & np_n-\log n\to c\in\mathbb{R},\\
1, & np_n-\log n\to +\infty.
\end{cases}
\]
The proof is organized into small lemmas.  The two main inputs are: (i) a detailed analysis of isolated vertices using alternating binomial sums, and (ii) a tree-counting argument showing that any other disconnected component is asymptotically negligible.  No branching-process heuristics are used.
\end{abstract}

\section{Statement of the theorem}

Throughout, $[n] := \{1,2,\dots,n\}$ and all logarithms are natural.  For $p\in[0,1]$, the random graph $G(n,p)$ is the graph on vertex set $[n]$ in which each edge is present independently with probability $p$.  Given a sequence $p_n\in[0,1]$, write
\[
\alpha_n := n p_n - \log n.
\]

\begin{theorem}[Connectivity threshold for $G(n,p)$]\label{thm:conn}
Let $G_n\sim G(n,p_n)$.  Then:
\begin{enumerate}
    \item If $\alpha_n\to -\infty$, then \(\PP(G_n\text{ is connected})\to 0\).
    \item If $\alpha_n\to c\in\mathbb{R}$, then \(\PP(G_n\text{ is connected})\to e^{-e^{-c}}\).
    \item If $\alpha_n\to +\infty$, then \(\PP(G_n\text{ is connected})\to 1\).
\end{enumerate}
\end{theorem}

The proof breaks into two independent pieces.

\begin{enumerate}
    \item Analyze the number of isolated vertices.  This already gives the correct limit for the probability of having \emph{no} isolated vertices.
    \item Show that, in the threshold window and above it, every disconnected graph with no isolated vertices has a component of size between $2$ and $n/2$, and the probability of such a component is $o(1)$.
\end{enumerate}


\section{Isolated vertices}

For each vertex $v\in[n]$, let
\[
I_v := \one_{\{v\text{ is isolated}\}},
\qquad
X_n := \sum_{v\in[n]} I_v.
\]
Thus $X_n$ is the number of isolated vertices of $G_n$.

\begin{lemma}[First two moments of the isolated-vertex count]
For every $n$ and every $p\in[0,1]$,
\[
\EE[X_n] = n(1-p)^{n-1}
\]
and
\[
\EE[X_n(X_n-1)] = n(n-1)(1-p)^{2n-3}.
\]
\end{lemma}

\begin{proof}
By linearity of expectation,
\[
\EE[X_n] = \sum_{v\in[n]} \PP(v\text{ is isolated}).
\]
A given vertex is isolated exactly when all $n-1$ incident edges are absent, which occurs with probability $(1-p)^{n-1}$.  This gives the first formula.

For the second formula, expand
\[
X_n(X_n-1) = \sum_{u\neq v} I_u I_v.
\]
Hence
\[
\EE[X_n(X_n-1)] = \sum_{u\neq v} \PP(u\text{ and }v\text{ are both isolated}).
\]
If $u\neq v$, then for both vertices to be isolated one must delete the $2n-4$ edges joining $u$ or $v$ to the other $n-2$ vertices, and also the edge $uv$.  In total this is $2n-3$ edges.  Therefore
\[
\PP(u\text{ and }v\text{ are both isolated})=(1-p)^{2n-3},
\]
which yields the second formula.
\end{proof}

\begin{proposition}[Subcritical regime: isolated vertices force disconnection]\label{prop:subcritical}
If $\alpha_n\to -\infty$, then
\[
\PP(X_n=0)\to 0.
\]
Consequently,
\[
\PP(G_n\text{ is connected})\to 0.
\]
\end{proposition}

\begin{proof}
Because $\alpha_n = np_n-\log n\to -\infty$, we have $np_n\le \log n$ for all large $n$, hence $p_n\le \log n/n\to 0$ and
\[
n p_n^2 \le \frac{\log^2 n}{n}\to 0.
\]
By the Taylor expansion $\log(1-x)=-x+O(x^2)$ for small $x$,
\[
(n-1)\log(1-p_n) = -(n-1)p_n + O(n p_n^2) = -n p_n + o(1).
\]
Therefore
\[
\EE[X_n]
= n(1-p_n)^{n-1}
= \exp\bigl(\log n - n p_n + o(1)\bigr)
= \exp\bigl(-\alpha_n + o(1)\bigr)\to \infty.
\]

Next, by the previous lemma,
\[
\frac{\EE[X_n(X_n-1)]}{\EE[X_n]^2}
=
\frac{n(n-1)(1-p_n)^{2n-3}}{n^2(1-p_n)^{2n-2}}
=
\frac{n-1}{n}(1-p_n)^{-1}
\to 1.
\]
Thus
\begin{align*}
\frac{\Var(X_n)}{\EE[X_n]^2}
&=
\frac{\EE[X_n]}{\EE[X_n]^2}
+
\frac{\EE[X_n(X_n-1)]}{\EE[X_n]^2}
-1 \\
&=
\frac{1}{\EE[X_n]}
+
\frac{\EE[X_n(X_n-1)]}{\EE[X_n]^2}
-1
\to 0.
\end{align*}
By Chebyshev's inequality,
\[
\PP(X_n=0)
\le
\PP\bigl(|X_n-\EE[X_n]|\ge \EE[X_n]\bigr)
\le
\frac{\Var(X_n)}{\EE[X_n]^2}
\to 0.
\]
Since every connected graph has no isolated vertex,
\[
\PP(G_n\text{ is connected}) \le \PP(X_n=0)\to 0.
\]
\end{proof}

The remaining cases require a more precise estimate on the probability of having no isolated vertices.

\begin{lemma}[Alternating binomial identity]
For integers $y\ge 1$ and $m\ge 0$,
\[
\sum_{r=0}^{m} (-1)^r \bbinom{y}{r} = (-1)^m \bbinom{y-1}{m}.
\]
In particular,
\[
\sum_{r=0}^{2m-1} (-1)^r \bbinom{y}{r} \le 0 \le \sum_{r=0}^{2m} (-1)^r \bbinom{y}{r}.
\]
\end{lemma}

\begin{proof}
Use Pascal's identity $\binom{y}{r}=\binom{y-1}{r}+\binom{y-1}{r-1}$:
\begin{align*}
\sum_{r=0}^{m} (-1)^r \bbinom{y}{r}
&=
\sum_{r=0}^{m} (-1)^r \bbinom{y-1}{r}
+
\sum_{r=1}^{m} (-1)^r \bbinom{y-1}{r-1} \\
&=
\sum_{r=0}^{m} (-1)^r \bbinom{y-1}{r}
-
\sum_{s=0}^{m-1} (-1)^s \bbinom{y-1}{s}
= (-1)^m \bbinom{y-1}{m}.
\end{align*}
The sign statement follows immediately because $\binom{y-1}{m}\ge 0$.
\end{proof}

\begin{lemma}[Bonferroni bounds for the event $\{Y=0\}$]
Let $Y$ be a nonnegative integer-valued random variable.  Then for every integer $m\ge 1$,
\[
\sum_{r=0}^{2m-1} (-1)^r \EE\!\left[\bbinom{Y}{r}\right]
\le
\PP(Y=0)
\le
\sum_{r=0}^{2m} (-1)^r \EE\!\left[\bbinom{Y}{r}\right].
\]
Here we use the convention $\binom{Y}{r}=0$ when $r>Y$.
\end{lemma}

\begin{proof}
Fix an outcome $\omega$ and write $y=Y(\omega)$.  If $y=0$, then $\binom{y}{0}=1$ and $\binom{y}{r}=0$ for $r\ge 1$, so both partial sums are exactly $1=\one_{\{y=0\}}$.

If $y\ge 1$, the previous lemma gives
\[
\sum_{r=0}^{2m-1} (-1)^r \bbinom{y}{r} \le 0 \le \sum_{r=0}^{2m} (-1)^r \bbinom{y}{r}.
\]
Since now $\one_{\{y=0\}}=0$, we have the pointwise bounds
\[
\sum_{r=0}^{2m-1} (-1)^r \bbinom{Y(\omega)}{r}
\le
\one_{\{Y(\omega)=0\}}
\le
\sum_{r=0}^{2m} (-1)^r \bbinom{Y(\omega)}{r}.
\]
Taking expectations proves the claim.
\end{proof}

\begin{proposition}[Critical window: the probability of having no isolated vertices]\label{prop:noiso}
Assume $\alpha_n\to c\in\mathbb{R}$.  Then
\[
\PP(X_n=0)\to e^{-e^{-c}}.
\]
\end{proposition}

\begin{proof}
For a fixed integer $r\ge 0$, let
\[
b_{n,r} := \EE\!\left[\bbinom{X_n}{r}\right].
\]
We first compute $b_{n,r}$ exactly.  The random variable $\binom{X_n}{r}$ counts $r$-element subsets of $[n]$ all of whose vertices are isolated.  Therefore, by symmetry,
\[
b_{n,r} = \bbinom{n}{r}\,\PP(\text{a fixed }r\text{-set consists entirely of isolated vertices}).
\]
If $R\subseteq[n]$ has size $r$, then every vertex of $R$ is isolated exactly when all edges with at least one endpoint in $R$ are absent.  The number of such edges is
\[
r(n-r)+\bbinom{r}{2} = rn - \frac{r(r+1)}{2}.
\]
Hence
\begin{equation}\label{eq:br-exact}
 b_{n,r} = \bbinom{n}{r}(1-p_n)^{r(n-r)+\binom{r}{2}}.
\end{equation}

Now fix $r$.  Since $\alpha_n\to c$, we have $p_n=(\log n + c + o(1))/n\to 0$ and
\[
n p_n^2 = p_n(np_n) = O\!\left(\frac{\log n}{n}\right)\to 0.
\]
Again using $\log(1-x)=-x+O(x^2)$,
\begin{align*}
\bigl(r(n-r)+\tbinom{r}{2}\bigr)\log(1-p_n)
&= -\bigl(rn+O(1)\bigr)p_n + O(n p_n^2) \\
&= -r n p_n + o(1)
= -r\log n - rc + o(1).
\end{align*}
Therefore
\[
(1-p_n)^{r(n-r)+\binom{r}{2}}
= n^{-r} e^{-rc + o(1)}.
\]
Also $\binom{n}{r}=n^r/r!+o(n^r)$.  Inserting this into \eqref{eq:br-exact} yields
\begin{equation}\label{eq:br-limit}
 b_{n,r} \to \frac{e^{-rc}}{r!}.
\end{equation}

Let $\lambda := e^{-c}$.  By the previous lemma, for every fixed $m\ge 1$,
\[
\sum_{r=0}^{2m-1} (-1)^r b_{n,r}
\le
\PP(X_n=0)
\le
\sum_{r=0}^{2m} (-1)^r b_{n,r}.
\]
Taking $n\to\infty$ and using \eqref{eq:br-limit} term by term,
\[
\sum_{r=0}^{2m-1} \frac{(-\lambda)^r}{r!}
\le
\liminf_{n\to\infty} \PP(X_n=0)
\le
\limsup_{n\to\infty} \PP(X_n=0)
\le
\sum_{r=0}^{2m} \frac{(-\lambda)^r}{r!}.
\]
Finally, send $m\to\infty$.  The even and odd partial sums of the exponential series converge to the same limit $e^{-\lambda}$.  Since $\lambda=e^{-c}$, we obtain
\[
\PP(X_n=0)\to e^{-\lambda}=e^{-e^{-c}}.
\]
\end{proof}

\begin{remark}
A standard refinement of the same computation shows that $X_n$ converges in distribution to a Poisson random variable of mean $e^{-c}$.  For the connectivity threshold, however, only the probability of the event $\{X_n=0\}$ is needed.
\end{remark}

\section{Nontrivial components are negligible}

Let $\mathcal{E}_n$ be the event that $G_n$ has a connected component of size between $2$ and $n/2$.

The next lemma gives a completely elementary upper bound for the probability that a prescribed vertex set forms a connected component.

\begin{lemma}[Probability that a fixed set is a connected component]
Let $S\subseteq[n]$ with $|S|=k\ge 2$.  Then
\[
\PP(S\text{ is a connected component of }G(n,p))
\le
k^{k-1} p^{k-1}(1-p)^{k(n-k)}.
\]
\end{lemma}

\begin{proof}
If $S$ spans a connected subgraph, then that subgraph contains a spanning tree on $S$.  Root such a tree at the smallest vertex of $S$.  Once the root is fixed, a rooted spanning tree is determined by the parent of each of the other $k-1$ vertices.  There are at most $k^{k-1}$ choices for these parents, so there are at most $k^{k-1}$ rooted spanning trees on $S$.

For any prescribed spanning tree on $S$, the probability that all its $k-1$ edges are present is $p^{k-1}$.  Independently, in order for $S$ to be a component, all $k(n-k)$ edges between $S$ and $[n]\setminus S$ must be absent, which has probability $(1-p)^{k(n-k)}$.  By the union bound over all rooted spanning trees,
\[
\PP(S\text{ is a connected component})
\le k^{k-1} p^{k-1}(1-p)^{k(n-k)}.
\]
\end{proof}

\begin{proposition}[No components of size $2$ through $n/2$ near the threshold]\label{prop:nosmall}
Assume that $\alpha_n\to c\in\mathbb{R}\cup\{+\infty\}$ and that
\[
p_n \le \frac{2\log n}{n}
\]
for all large $n$.  Then
\[
\PP(\mathcal{E}_n)\to 0.
\]
\end{proposition}

\begin{proof}
Write $p=p_n$ and $\alpha=np-\log n$.  Since $\alpha_n\to c\in\mathbb{R}\cup\{+\infty\}$, we have $np=\log n+O(1)$ or larger; in particular $np\to\infty$.

For each $k$, let $A_{n,k}$ be the event that $G_n$ has a component of size exactly $k$.  By the union bound over $k$-subsets and the previous lemma,
\begin{equation}\label{eq:ank-raw}
\PP(A_{n,k})
\le
\bbinom{n}{k} k^{k-1} p^{k-1}(1-p)^{k(n-k)}.
\end{equation}
Using $\binom{n}{k}\le (en/k)^k$ and $(1-p)^m\le e^{-pm}$,
\begin{equation}\label{eq:ank-bound}
\PP(A_{n,k})
\le
\frac{k}{p}\bigl(e n p \, e^{-p(n-k)}\bigr)^k.
\end{equation}
We now sum this bound over $2\le k\le n/2$.

\smallskip
\textbf{Step 1: small components.}
Let
\[
K_n := \left\lfloor \frac{n}{\log^2 n}\right\rfloor.
\]
Suppose $2\le k\le K_n$.  Since $p\le 2\log n/n$,
\[
pk \le \frac{2\log n}{n}\cdot \frac{n}{\log^2 n} = \frac{2}{\log n},
\]
so $e^{pk}\le 2$ for all large $n$.  Therefore
\[
e^{-p(n-k)} = e^{-pn}e^{pk} \le 2e^{-pn}
\]
and \eqref{eq:ank-bound} gives
\[
\PP(A_{n,k}) \le \frac{k}{p} \rho_n^k,
\qquad
\rho_n := 2e n p e^{-pn} = 2e (np)e^{-np}.
\]
Because $np\to\infty$, we have $\rho_n\to 0$.  Hence, for large $n$, $\rho_n\le 1/2$, and then
\begin{align*}
\sum_{k=2}^{K_n} \PP(A_{n,k})
&\le \frac{K_n}{p}\sum_{k=2}^{\infty} \rho_n^k
\le \frac{2K_n\rho_n^2}{p} \\
&= 8e^2 K_n n^2 p e^{-2np}
= 8e^2 K_n p e^{-2\alpha}.
\end{align*}
Since $K_n p\le 2/\log n$ and $e^{-2\alpha}$ is bounded if $\alpha_n\to c\in\mathbb{R}$ and tends to $0$ if $\alpha_n\to +\infty$, the right-hand side tends to $0$.

\smallskip
\textbf{Step 2: larger components.}
Now suppose $K_n\le k\le n/2$.  Then $n-k\ge n/2$, so \eqref{eq:ank-bound} implies
\[
\PP(A_{n,k})
\le
\frac{k}{p} q_n^k,
\qquad
q_n := e n p e^{-pn/2} = e(np)e^{-np/2}.
\]
Again $np\to\infty$, so $q_n\to 0$.  Also $k\le 2^k$ for all $k\ge 1$, hence
\[
\PP(A_{n,k}) \le \frac{(2q_n)^k}{p}.
\]
For all large $n$, we have $2q_n\le 1/2$, and therefore
\[
\sum_{k=K_n}^{\lfloor n/2\rfloor} \PP(A_{n,k})
\le
\frac{1}{p}\sum_{k=K_n}^{\infty} (2q_n)^k
\le
\frac{2^{1-K_n}}{p}.
\]
Because $\alpha_n$ is bounded below eventually, we have $np = \log n + \alpha_n \ge \tfrac12\log n$ for all large $n$, so $1/p = O(n/\log n)$.  Since $K_n\to\infty$, the bound $2^{1-K_n}/p$ tends to $0$.

Combining the two ranges and summing over $k$ proves
\[
\PP(\mathcal{E}_n)
\le
\sum_{k=2}^{\lfloor n/2\rfloor} \PP(A_{n,k})
\to 0.
\]
\end{proof}

\section{Monotonicity in \texorpdfstring{$p$}{p}}

\begin{lemma}[Connectivity is increasing in the edge parameter]\label{lem:monotone}
If $0\le p\le q\le 1$, then
\[
\PP\bigl(G(n,p)\text{ is connected}\bigr)
\le
\PP\bigl(G(n,q)\text{ is connected}\bigr).
\]
\end{lemma}

\begin{proof}
Assign to each potential edge $e$ an independent random variable $U_e\sim \mathrm{Unif}[0,1]$.  Construct two random graphs on the same probability space by declaring
\[
e\in G(n,p) \iff U_e\le p,
\qquad
e\in G(n,q) \iff U_e\le q.
\]
Because $p\le q$, every edge present in $G(n,p)$ is also present in $G(n,q)$.  Thus, under this coupling,
\[
G(n,p)\subseteq G(n,q)
\]
as spanning subgraphs on the same vertex set.  Connectivity is an increasing graph property, so
\[
\{G(n,p)\text{ is connected}\}
\subseteq
\{G(n,q)\text{ is connected}\}.
\]
Taking probabilities gives the claim.
\end{proof}

\section{Proof of the main theorem}

\begin{proof}[Proof of Theorem~\ref{thm:conn}]
We treat the three cases separately.

\smallskip
\textbf{Case 1: $\alpha_n\to -\infty$.}
By Proposition~\ref{prop:subcritical}, the probability that there are no isolated vertices tends to $0$.  Hence the graph is connected with probability tending to $0$.

\smallskip
\textbf{Case 2: $\alpha_n\to c\in\mathbb{R}$.}
Since $np_n=\log n + O(1)$, we have $p_n\le 2\log n/n$ for all large $n$.  Proposition~\ref{prop:nosmall} gives
\[
\PP(\mathcal{E}_n)\to 0.
\]
A disconnected graph either has an isolated vertex or has a component of size between $2$ and $n/2$.  Therefore
\[
0
\le
\PP(X_n=0) - \PP(G_n\text{ is connected})
\le
\PP(\mathcal{E}_n).
\]
By Proposition~\ref{prop:noiso},
\[
\PP(X_n=0)\to e^{-e^{-c}},
\]
so the squeeze theorem gives
\[
\PP(G_n\text{ is connected})\to e^{-e^{-c}}.
\]

\smallskip
\textbf{Case 3: $\alpha_n\to +\infty$.}
Define
\[
q_n := \min\!\left\{p_n,\frac{2\log n}{n}\right\}
\]
and let $H_n\sim G(n,q_n)$.  Then $q_n\le 2\log n/n$ for all $n$, and
\[
nq_n - \log n \to +\infty.
\]
Indeed, if $q_n=p_n$ then this is exactly the hypothesis; otherwise $q_n=2\log n/n$ and $nq_n-\log n = \log n\to\infty$.

By the union bound,
\[
\PP(H_n\text{ has an isolated vertex})
\le
n(1-q_n)^{n-1}
\le
\exp\bigl(\log n - q_n(n-1)\bigr)
\to 0,
\]
because $nq_n-\log n\to\infty$.

Also Proposition~\ref{prop:nosmall} applies to $H_n$, so
\[
\PP(H_n\text{ has a component of size between }2\text{ and }n/2)\to 0.
\]
Hence
\[
\PP(H_n\text{ is disconnected})\to 0,
\]
that is,
\[
\PP(H_n\text{ is connected})\to 1.
\]
Finally, by Lemma~\ref{lem:monotone},
\[
\PP(G_n\text{ is connected})
\ge
\PP(H_n\text{ is connected})
\to 1.
\]
This proves the supercritical case.
\end{proof}

\end{document}
